% ============================================================================
% SEC03.TEX
% Section 1.3: Game Design Document (GDD)
% ============================================================================

\section{Game Design Document (GDD)}

\begin{learningoutcomes}
\item Memahami tujuan dan fungsi Game Design Document
\item Mengetahui struktur dan komponen GDD
\item Mampu menyusun GDD untuk konsep game
\item Memahami pentingnya dokumentasi dalam pengembangan game
\end{learningoutcomes}

\subsection{Pengertian Game Design Document}

\begin{definisi}[Game Design Document]
Game Design Document (GDD) adalah dokumen komprehensif yang menjelaskan konsep, mekanik, desain, dan visi keseluruhan dari sebuah game \cite{Schell2019}.
\end{definisi}

GDD berfungsi sebagai:
\begin{itemize}
\item Blueprint untuk pengembangan game
\item Komunikasi antara anggota tim
\item Referensi selama proses development
\item Dokumentasi untuk investor atau publisher
\end{itemize}

\subsection{Struktur GDD}

\subsubsection{1. Executive Summary}
Ringkasan singkat tentang game:
\begin{itemize}
\item Judul game
\item Genre
\item Platform target
\item Target audience
\item Unique selling point
\end{itemize}

\subsubsection{2. Game Overview}
Deskripsi umum game:
\begin{itemize}
\item Konsep inti
\item Tema dan setting
\item Tone dan mood
\item Visual style
\end{itemize}

\subsubsection{3. Gameplay Mechanics}
Detail mekanik gameplay:
\begin{itemize}
\item Core mechanics
\item Control scheme
\item Progression system
\item Win/lose conditions
\end{itemize}

\subsubsection{4. Story and Narrative}
Elemen naratif:
\begin{itemize}
\item Plot overview
\item Karakter utama
\item World building
\item Dialog system (jika ada)
\end{itemize}

\subsubsection{5. Art and Visual Design}
Arah visual:
\begin{itemize}
\item Art style
\item Color palette
\item Character design
\item Environment design
\end{itemize}

\subsubsection{6. Audio Design}
Elemen audio:
\begin{itemize}
\item Music style
\item Sound effects
\item Voice acting (jika ada)
\item Audio implementation
\end{itemize}

\subsubsection{7. Technical Specifications}
Spesifikasi teknis:
\begin{itemize}
\item Engine yang digunakan
\item Platform requirements
\item Performance targets
\item Technical features
\end{itemize}

\subsubsection{8. Development Timeline}
Jadwal pengembangan:
\begin{itemize}
\item Milestones
\item Resource requirements
\item Risk assessment
\end{itemize}

\begin{catatan}
GDD adalah dokumen hidup yang dapat berubah selama proses development. Penting untuk menjaga GDD tetap update dengan perubahan desain.
\end{catatan}

\subsection{Best Practices dalam Menulis GDD}

\begin{enumerate}
\item \textbf{Jelas dan Spesifik}: Hindari ambiguitas
\item \textbf{Visual}: Gunakan diagram, mockup, dan referensi visual
\item \textbf{Organized}: Struktur yang mudah diikuti
\item \textbf{Accessible}: Mudah dipahami oleh semua anggota tim
\item \textbf{Version Control}: Kelola perubahan dengan baik
\end{enumerate}

\subsection{Contoh Template GDD}

\begin{contoh}[Template GDD Sederhana]
\textbf{1. Game Title}: [Nama Game]

\textbf{2. Genre}: [Genre Utama]

\textbf{3. Platform}: [PC/Mobile/Console]

\textbf{4. Target Audience}: [Demografi]

\textbf{5. Core Gameplay Loop}: [Deskripsi singkat gameplay utama]

\textbf{6. Key Features}: 
\begin{itemize}
\item Feature 1
\item Feature 2
\item Feature 3
\end{itemize}

\textbf{7. Art Style}: [Deskripsi visual style]

\textbf{8. Development Timeline}: [Estimasi waktu]
\end{contoh}
